\documentclass[11pt,aps,preprint,nofootinbib,superscriptaddress,tightenlines]{revtex4-2}
\usepackage[T1]{fontenc}
\usepackage{amsmath,amssymb,mathtools,bm}
\usepackage{xcolor}
\usepackage{booktabs}
\usepackage{graphicx}
\usepackage{dsfont}
\newcommand{\one}{\mathds{1}}
\newcommand{\avg}[1]{\langle #1\rangle}
\definecolor{darkblue}{RGB}{20,60,140}
\usepackage[colorlinks=true,linkcolor=darkblue,citecolor=darkblue,urlcolor=darkblue]{hyperref}

\def\st{\begin{equation}}
\def\stp{\end{equation}}
\def\Eq#1{eq.~(\ref{#1})}
\def\Sect#1{Sect.~\ref{#1}}
\def\Fig#1{Fig.~\ref{#1}}
\let\oldparagraph\paragraph
\renewcommand{\paragraph}{\par\bigskip\oldparagraph}

\begin{document}
\sloppy

\title{The inputs of a run}
\author{D.~Teaney}
\affiliation{Department of Physics and Astronomy, Stony Brook University,
Stony Brook, New York 11794, USA}
\date{\today}

\begin{abstract}
A short review of what a run of \texttt{examples/equilibrium.jl} (and of \texttt{examples/hmc\_positions.jl}) is given,
and how the inputs become the variables of the particles.  The physics is a thermal state $(T,\bm v)$ of a fluid with a given equation of
state and viscosity in a periodic box; the numerics are the particle number $N$ and the smoothing length $h$.  These combine into the
entropy per particle $s_a$, the parameter that controls the fluctuations of a particle and the size of the renormalization.  The note is
written for the defaults after the merge of branch \texttt{split} (\texttt{SPH\_RENORM=split}).
\end{abstract}

\maketitle

%=============================================================================
\section{The inputs in one view}
\label{sec:overview}

A run is specified by three groups of inputs.
\begin{enumerate}
\item \emph{The physics}: the equation of state (speed of sound $c_s^2$, temperature scale $T_{\rm scale}$), the transport coefficients
  ($\eta/s$, $\zeta/s$), and the thermal state, a rest-frame temperature $T_{\rm rel}$ and a velocity $\bm v$.
\item \emph{The box}: a periodic square of side $L$, volume $L^2$.
\item \emph{The discretization}: the smoothing length $h$ and the particle spacing $a$, given as $h/a$; together with $L$ they fix the
  particle number $N$ and the number density $n=N/L^2$.
\end{enumerate}
Everything else (the integrators, the measure, the renormalization, the start) is a numerical setting with a default (\Sect{sec:numerics}).
The inputs meet in one number, the entropy per particle,
\st
s_a=\frac{S}{n},
\label{eq:sa}
\stp
with $S$ the entropy density of the state; $s_a$ is what the particles carry.

\paragraph{What the drivers take today.}  The drivers do not take the temperature: they take the \emph{rest-frame entropy density} $s_{\rm rel,0}$
(\texttt{SR\_SREL0}, formerly \texttt{SR\_S0}, retired 2026-11-23), and derive the temperature from it through the equation of state,
\st
T_{\rm rel,0}=T_{\rm rel}\big|_{s_{\rm rel}=s_{\rm rel,0}}\qquad\text{(uniform fluid at rest)} .
\label{eq:Trel0}
\stp
The run is then at a definite $(T_{\rm rel,0},\bm v)$ (\Sect{sec:start}), but $T_{\rm rel,0}$ is derived, not given.  A proposal to take
the temperature directly is in \Sect{sec:proposal}.

%=============================================================================
\section{The physics}
\label{sec:physics}

\paragraph{The equation of state} (\texttt{ConstantCs2EOS}, \texttt{src/sph.jl}) has a constant speed of sound:
\st
p=c_s^2\,e,\qquad T_{\rm rel}=T_{\rm scale}\,e^{c_s^2/(1+c_s^2)},\qquad s_{\rm rel}=\frac{e+p}{T_{\rm rel}} ,
\label{eq:eos}
\stp
with $e$, $p$, $s_{\rm rel}$ the rest-frame energy density, pressure and entropy density ($S=\gamma s_{\rm rel}$ is the entropy density in the lab, the density frame; $s_{\rm rel,0}$ is the value of a run, \texttt{SR\_SREL0}).  Eliminating $e$,
\st
s_{\rm rel}=(1+c_s^2)\,T_{\rm scale}^{-(1+c_s^2)/c_s^2}\;T_{\rm rel}^{1/c_s^2} .
\label{eq:sofT}
\stp
The drivers use the conformal fluid in $d=2$, $c_s^2=1/d=1/2$ and $T_{\rm scale}=1$ (\texttt{conformal\_eos(2)}), for which
\st
e=T_{\rm rel}^3,\qquad p=\tfrac12T_{\rm rel}^3,\qquad s_{\rm rel}=\tfrac32T_{\rm rel}^2 .
\label{eq:conformal}
\stp
\texttt{SR\_CS2} selects another constant $c_s^2$ (e.g.\ $1/3$, $e=T_{\rm rel}^4$, $s_{\rm rel}=\tfrac43T_{\rm rel}^3$), still with $T_{\rm scale}=1$.

\paragraph{The temperatures.}  The particles move with the bare temperature $T=\partial H/\partial s$ and the rest-frame temperature is
\st
T_{\rm rel}=\gamma\,T ,\qquad \gamma=(1-v^2)^{-1/2};
\stp
$T_{\rm rel}$ enters only with the viscosities.

\paragraph{The transport coefficients} are given as $\eta/s$ (\texttt{SR\_ETAS}, default $0.2$) and $\zeta/s$ (\texttt{SR\_ZETAS},
default $0$).  The kinematic shear viscosity of the state is
\st
\nu_0=\frac{\eta}{e+p}=\frac{\eta/s}{T_{\rm rel}} .
\label{eq:nu0}
\stp
The viscous step uses the effective bulk viscosity
\st
\zeta_{\rm eff}=\max\Big(\zeta+2\eta\Big(\frac12-\frac1d\Big),\;\epsilon\,\eta\Big),\qquad \epsilon=0.01\ (\texttt{SPH\_ZETA\_FLOOR}),
\stp
which at $d=2$, $\zeta=0$ is the floor $0.01\,\eta$.

\paragraph{The state.}  The fluid is uniform, with rest-frame entropy density $s_{\rm rel,0}$ and velocity $\bm v=v\hat{\bm x}$, $v=$
\texttt{SR\_BOOST} (default $0$).  In the lab (the density frame) its entropy and momentum densities are
\st
S=\gamma s_{\rm rel,0},\qquad \bm M=w_0\gamma^2\bm v,\qquad w_0=e+p\big|_{s_{\rm rel,0}},
\label{eq:lab}
\stp
with $\bm M$ the density-frame momentum density.

%=============================================================================
\section{The box and the discretization}
\label{sec:discretization}

\paragraph{The numbers given.}  The command line of \texttt{equilibrium.jl} is
\begin{center}
\texttt{julia -t auto --project examples/equilibrium.jl L h/a t\_end nreal [rec\_every] [tag]}
\end{center}
The smoothing length is not on the command line: it is \texttt{SR\_H} (default $h=0.05$).  The kernel is the Wendland C2 kernel
(\texttt{SPH\_KERNEL=wc2}) with support $2h$.

\paragraph{The lattice} (\texttt{SimParams}, \texttt{src/sph.jl}).  The spacing is
\st
a=\frac{h}{h/a},
\stp
and the particles are placed on a hexagonal lattice: $N_x$ columns of spacing $a$ and $N_y$ rows of spacing $a\sqrt3/2$,
\st
N_x=\mathrm{round}\Big(\frac La\Big),\qquad N_y=2\,\mathrm{round}\Big(\frac{L}{a\sqrt3}\Big),\qquad N=N_xN_y ,
\stp
($N_y$ even for the periodic offset rows).  The box is snapped to the lattice, $L_x=N_xa$, $L_y=N_ya\sqrt3/2$, so it is square only
approximately ($L=0.6$, $h/a=2$: $L_x=0.6$, $L_y=0.606$).  The volume per particle and the number density are
\st
V=\frac{\sqrt3}{2}\,a^2=\frac1n,\qquad n=\frac{N}{L_xL_y}=\frac{2}{\sqrt3}\,\frac{(h/a)^2}{h^2} .
\label{eq:n}
\stp
So $L$, $h$ and $h/a$ are the three independent numbers; $N$ and $n$ follow.  The number of particles within the kernel's support is
$\pi(2h)^2n$ ($58$ at $h/a=2$).

\paragraph{The particle variables} (\texttt{setup\_ic}, \texttt{src/setup.jl}).  Each particle receives the densities of \Eq{eq:lab}
times its lattice volume,
\st
s_a=\gamma s_{\rm rel,0}\,V=\frac{\gamma s_{\rm rel,0}}{n},\qquad \bm m_a=\bm M\,V .
\label{eq:particles}
\stp
The rest-frame value $s_{\rm rel,0}V$ is the $s_a$ quoted for a run.  With \texttt{SPH\_RENORM=split} the particle variable $s_a$ is the label
$u$, the entropy of the target state of the particle (\texttt{docs/step5\_plan.tex}); it is set up by \Eq{eq:particles} in the same way.

%=============================================================================
\section{\texorpdfstring{The start: the state is $(T_{\rm rel,0},\bm v)$}{The start}}
\label{sec:start}

The set-up of \Eq{eq:particles} is a lattice with no fluctuations.  The start replaces it by a sample of the equilibrium
(\texttt{docs/scheme.tex}, Sect.~``The start''): hybrid Monte Carlo of the positions, energies and momenta (\texttt{SR\_POSITIONS=hmc}),
then an equilibrium draw of the energies and momenta at the sampled positions (\texttt{SR\_START=eq}).  Both are at the multipliers
of the requested state (\texttt{SR\_BATH=nominal}, the default since commit \texttt{3262b00}, 2026-11-23),
\st
\beta^0=\frac1T=\frac{\gamma}{T_{\rm rel,0}},\qquad \bm\beta=\frac{\bm v}{T}=\frac{\gamma\bm v}{T_{\rm rel,0}} ,
\label{eq:bath}
\stp
with $T_{\rm rel,0}$ the temperature of the uniform rest-frame state at $s_{\rm rel,0}$, \Eq{eq:Trel0}.  The draw then shifts the energies and momenta
so that the totals are exactly those of the state it starts from.  The fluid of the run is at this $(T_{\rm rel,0},\bm v)$ for any
Hamiltonian, tree level or renormalized.  (\texttt{SR\_BATH=totals}, before 2026-11-23, solved for $\beta^\mu$ from the set-up totals,
which at low $s_a$ in a moving fluid could be far from the requested state; \texttt{docs/split\_dynamics.tex}.)

In short: the input is $s_{\rm rel,0}$; the state is $(T_{\rm rel,0},\bm v)$ with
\st
T_{\rm rel,0}=\Big[\frac{s_{\rm rel,0}\,T_{\rm scale}^{(1+c_s^2)/c_s^2}}{1+c_s^2}\Big]^{c_s^2}
\quad\Big(=\sqrt{2s_{\rm rel,0}/3}\ \text{ at }c_s^2=\tfrac12,\ T_{\rm scale}=1\Big),
\label{eq:TofS}
\stp
and the particles carry $s_a=\gamma s_{\rm rel,0}/n$ at the start, which then fluctuates.

%=============================================================================
\section{Dimensionless parameters and derived numbers}
\label{sec:derived}

\paragraph{Scale invariance.}  The equation of state \Eq{eq:eos} has no scale besides $T_{\rm scale}$, so under $\bm x\to\lambda\bm x$,
$t\to\lambda t$, $T_{\rm rel}\to T_{\rm rel}/\lambda$ the run is unchanged.  A run is therefore fixed by the dimensionless numbers
\st
T_{\rm rel,0}\,h\ \ (\text{equivalently } s_{\rm rel,0}h^2),\qquad nh^2\ \ (\text{equivalently } h/a),\qquad \frac Lh,\qquad v,\qquad \frac\eta s,\quad \frac\zeta s,\quad c_s^2 ,
\label{eq:dimless}
\stp
and the entropy per particle is the ratio of the first two,
\st
s_a=\frac{s_{\rm rel,0}h^2}{nh^2}\qquad(\text{at rest}).
\label{eq:saratio}
\stp
The bare viscosity in units of $h$ is $\nu_0/h=(\eta/s)/(T_{\rm rel,0}h)$.

\paragraph{The standard runs.}  $L=0.6$, $h=0.05$, $h/a=2$: $a=0.025$, $N=672$, $n=1848$, $nh^2=4.62$, the smallest box mode
$kh=2\pi h/L=0.52$.
\begin{center}
\begin{tabular}{cccccccc}
\toprule
$s_{\rm rel,0}$ & $s_{\rm rel,0}h^2$ & $s_a$ & $c_s^2$ & $c_s$ & $T_{\rm rel,0}$ & $\nu_0=(\eta/s)/T_{\rm rel,0}$ & $\nu_0/h$ \\ \midrule
$2.956\times10^4$ & $73.9$ & $16$ & $1/2$ & $0.707$ & $140.4$ & $1.42\times10^{-3}$ & $0.028$ \\
                  &        &      & $1/3$ & $0.577$ & $28.1$  & $7.12\times10^{-3}$ & $0.142$ \\
$9.977\times10^4$ & $249$  & $54$ & $1/2$ & $0.707$ & $257.9$ & $7.75\times10^{-4}$ & $0.016$ \\
                  &        &      & $1/3$ & $0.577$ & $42.1$  & $4.75\times10^{-3}$ & $0.095$ \\
\bottomrule
\end{tabular}
\end{center}
$\eta/s=0.2$.  The default $s_{\rm rel,0}=10^6$ gives $s_a=541$ in the same box.

\paragraph{Consequence for the damping.}  At the same $s_{\rm rel,0}$ the bare $\nu_0/h$ at $c_s^2=1/2$, $\eta/s=0.2$ is five times smaller than
at $c_s^2=1/3$, and the measured shear damping at $kh=0.52$ is then $3$--$4$ times hydrodynamics, a fluctuation correction
$\propto1/(s\,\nu_0)$ to the viscosity; $\eta/s=1$ (or $c_s^2=1/3$) brings it back to $1.1$ times hydrodynamics
(\texttt{docs/split\_dynamics.tex}, campaigns \texttt{damping\_scan}, \texttt{cs13}).

%=============================================================================
\section{Refinement and the renormalization}
\label{sec:refinement}

At fixed physics (fixed $T_{\rm rel,0}$, $\bm v$, $\eta/s$, $L$) and fixed coarse-graining scale $h$, a refinement increases the number
density $n$ at fixed $s_{\rm rel,0}h^2$:
\st
s_{\rm rel,0}h^2=s_a\,nh^2=\text{fixed},\qquad
\frac ha=\Big(\frac{\sqrt3}{2}\,nh^2\Big)^{1/2}\propto\sqrt n,\qquad
s_a=\frac{s_{\rm rel,0}}{n}\propto\frac1n .
\label{eq:refine}
\stp
(The same holds for a change of the scale $h$ at fixed $n$: what matters is the pair $s_{\rm rel,0}h^2$, $nh^2$ of \Eq{eq:saratio}.)
The particles become lighter: each fluctuates with a relative size $1/\sqrt{s_a}$.  At tree level the canonical ensemble of the particle
Hamiltonian has the target thermodynamics plus corrections of order $T$ per particle, i.e.\ of relative size
\st
\frac{1}{s_a}=\frac{nh^2}{s_{\rm rel,0}h^2}\ \propto\ n ,
\label{eq:oneloop}
\stp
which grow linearly with the refinement (at $s_a=16$ the tree-level equal-time spectra are $20$--$40\%$ above the susceptibilities of
the equation of state).  A continuum limit at fixed physics therefore needs counterterms that remove these terms particle by particle:
this is the renormalization, \texttt{SPH\_RENORM=split} (\texttt{docs/step5\_plan.tex}), with which the spectra agree with the target to a
few per cent at $s_a=16$.

%=============================================================================
\section{The numerical settings}
\label{sec:numerics}

All are environment variables, read through \texttt{setting} (\texttt{src/settings.jl}) and written to each run's record (\texttt{.toml}).
\begin{center}\small
\begin{tabular}{lll}
\toprule
setting & default & alternatives \\ \midrule
\texttt{SPH\_HAMILTONIAN}, \texttt{SPH\_VOLUME} & \texttt{energy}, \texttt{number} & \texttt{smoothed}, \texttt{entropy} \\
\texttt{SPH\_IDEAL} & \texttt{generalized} & \texttt{particle}, \texttt{grid} \\
\texttt{SPH\_DISSIPATION} & \texttt{moment} & \texttt{common}, \texttt{particle} \\
\texttt{SPH\_HASTINGS} & on where it applies & \texttt{0} \\
\texttt{SPH\_MEASURE} & \texttt{liouville} & \texttt{einstein} \\
\texttt{SPH\_GAMMA\_MAX} & $100$ & \\
\texttt{SPH\_ZETA\_FLOOR} & $0.01$ & $0$ (no floor) \\
\texttt{SPH\_RENORM} & \texttt{split}$^\ast$ & \texttt{legendre}, \texttt{1} (label only), \texttt{0} (tree level) \\
\texttt{SPH\_RENORM\_VBAR} & $L_xL_y/N$ & (enters only for $c_s^2\neq1/2$) \\
\texttt{SPH\_KERNEL} & \texttt{wc2} & \texttt{wc4}, \texttt{wc6} \\
\texttt{SR\_POSITIONS} & \texttt{hmc} & \texttt{lattice}, \texttt{poisson} \\
\texttt{SR\_HMC\_N}, \texttt{\_TAU}, \texttt{\_TAUJIT}, \texttt{\_DTFAC} & $200$, $L/c_s$, $0.5$, $0.35$ & \\
\texttt{SR\_START} & \texttt{eq} & \texttt{cold}, \texttt{thermal} \\
\texttt{SR\_BATH} & \texttt{nominal} & \texttt{totals} \\
\texttt{SR\_DTFAC} & $1$ & \\
\bottomrule
\end{tabular}
\end{center}
$^\ast$Default from the merge of branch \texttt{split}; on the branch today the default is the tree level.  The time step is
\st
\Delta t=\tfrac12\,\texttt{SR\_DTFAC}\,\min\Big(\frac{0.2\,h}{c_s},\ \frac{0.2\,h^2\,T_{\rm rel,min}}{\eta/s+\zeta/s},\ \frac{2}{\Gamma_{\rm bound}}\Big),
\stp
the acoustic and diffusive limits and the bound $\Gamma_{\rm bound}$ on the damping rate of the viscous step (\texttt{compute\_dt},
\texttt{src/ideal.jl}; the diffusive term written for $d=2$).  Hybrid Monte Carlo uses $\delta t=\texttt{SR\_HMC\_DTFAC}\times0.2h/c_s$.

%=============================================================================
\section{All inputs}
\label{sec:table}

\begin{center}\small
\begin{tabular}{llll}
\toprule
input & symbol & setting & default \\ \midrule
box side & $L$ & argument 1 & --- \\
resolution & $h/a$ & argument 2 & --- \\
smoothing length & $h$ & \texttt{SR\_H} & $0.05$ \\
rest-frame entropy density & $s_{\rm rel,0}$ & \texttt{SR\_SREL0} & $10^6$ \\
velocity (along $x$) & $v$ & \texttt{SR\_BOOST} & $0$ \\
speed of sound & $c_s^2$ & \texttt{SR\_CS2} & $0$, i.e.\ $1/d=1/2$ \\
temperature scale & $T_{\rm scale}$ & (fixed) & $1$ \\
shear viscosity & $\eta/s$ & \texttt{SR\_ETAS} & $0.2$ \\
bulk viscosity & $\zeta/s$ & \texttt{SR\_ZETAS} & $0$ \\
run length & $t_{\rm end}$ & argument 3 & --- \\
realizations & & argument 4 & --- \\
record every (steps) & & argument 5 & $2$ \\
tag & & argument 6 & \texttt{L<L>\_ha<h/a>} \\
largest recorded $kh$ & & \texttt{SR\_KHMAX} & $1.5$ \\
seeds & & \texttt{SR\_SEED0}, \texttt{SR\_SEEDSTRIDE}, \texttt{SR\_SEEDOFFSET} & $0$, $1$, $0$ \\
\midrule
\multicolumn{4}{l}{derived: $a$, $N$, $n$, $s_a=\gamma s_{\rm rel,0}/n$, $T_{\rm rel,0}$ (\Eq{eq:TofS}), $\nu_0$, $\Delta t$} \\
\bottomrule
\end{tabular}
\end{center}

%=============================================================================
\section{Remarks from the code}
\label{sec:remarks}

\paragraph{A proposal: the temperature as input.}
\label{sec:proposal}
A setting \texttt{SR\_T} $=T_{\rm rel,0}$ could replace \texttt{SR\_SREL0}: the driver would set $s_{\rm rel,0}$ from \Eq{eq:sofT} and proceed
unchanged.  The run is already at $(T_{\rm rel,0},\bm v)$ through \Eq{eq:bath}; only the way it is named would change, and the inputs
would then read as the physics does: $T$, $v$, the equation of state, $\eta/s$; then $L$, $N$ (or $h/a$) and $h$.  Not implemented.

\paragraph{Points found in reading the code.}
\begin{itemize}
\item In \texttt{equilibrium.jl} the default HMC trajectory length is \texttt{LBOX*sqrt(3)}, which is $L/c_s$ only at $c_s^2=1/3$; at
  $c_s^2=1/2$ it is $\sqrt{3/2}\,L/c_s$.  The comments (and \texttt{split\_dynamics.tex}) say $L/c_s$.  With the random lengths
  (\texttt{SR\_HMC\_TAUJIT=0.5}) this is harmless, but the statement or the code should be made to agree.
\item \texttt{hmc\_positions.jl} (arguments \texttt{L h/a nupdates nsteps [tag]}) is narrower than \texttt{equilibrium.jl}: $h=0.05$,
  $\eta/s=0.2$ and $c_s^2=1/2$ are fixed in the code (no \texttt{SR\_H}, \texttt{SR\_ETAS}, \texttt{SR\_CS2}), and its defaults differ:
  \texttt{SR\_HMC\_DTFAC} $0.5$, \texttt{SR\_HMC\_TAUJIT} $0$, \texttt{SR\_SEED0} $1$, $\tau=L/c_s$ (computed from $c_s$).
\item The box is not exactly square (\Sect{sec:discretization}); $n$ and the box modes use $L_x$, $L_y$ as snapped.
\item The quoted $s_a$ is the set-up value $\gamma s_{\rm rel,0}V$; with \texttt{SR\_POSITIONS=hmc} the volumes $V_a=1/n_a$ and the $s_a$
  fluctuate from the start, and their mean need not equal it.
\end{itemize}

\end{document}
